\section{Discussion}
\label{sec:discussion}

\paragraph{What decision models buy.} The clearest benefit is the absence of a labelled training set: even with label names only,
no compact classifier we tested came close with 5 or 20 examples per class. The second benefit is operational: the Base Labs models gave the same answer
for any option order, had well-calibrated probabilities and therefore supported a usable ``accept automatically or send to a human''
policy on Banking77, with an observed error close to, though not always below, the target. Neither benefit depends on the decision models being far better than a well-trained small model when thousands of
labels are available; in that regime the remaining gap was 0--2.4 points on two of the three datasets (5.5 on MASSIVE-it),
and on Banking77 a different training recipe left no resolved difference.

\paragraph{Where the advantage comes from.} Comparing Elephant-1 with Qwen3-8B suggested a large effect of decision-specific
post-training; compared with its actual base, Qwen3.5-9B, the observed gap is much smaller. This is not a causal decomposition (the two
also differ in attention mask, objective, precision and serving), but it shows that an older generalist is the wrong baseline.
For a practitioner this has a practical consequence: a recent open generalist, prompted with good category descriptions and scored
on option letters, is already a strong zero-shot classifier on a task like Ottavo, although it is slower in our unoptimised
implementation and more sensitive to option order.

\paragraph{When a small classifier is the better choice.} Our data support a compact classifier when (i) thousands of labelled
examples exist or can be produced at an acceptable cost, (ii) the taxonomy is stable, so labels do not need to be redone, and (iii)
volume, latency or data-residency constraints matter more than one or two points of macro-F1. Below about 20 examples per class,
SetFit was the best compact option among those we evaluated; decoders with a new classification head collapsed with both recipes we tried.

\paragraph{Descriptions, prompts and evaluation practice.} Category descriptions and in-context examples are not universally helpful:
their effect varies with the model and the dataset. Descriptions written from labelled examples also spend labels, which a fair
comparison with trained classifiers has to count. Practitioners should measure both on their own data, using the same messages for the comparison, and
prefer retrieved examples to long lists of examples when a prompt-based model benefits from them. More generally, evaluations that report
a single run on a single test sample can mislead in both directions. The exploratory result that motivated this study (a small model
beating a decision model on Banking77) looked refuted under our frozen protocol, which scored the small model 2.6 points lower; the
exploratory recipe, re-run with four seeds, reproduced the original score, with no resolved difference from the decision model. A frozen protocol protects
against choosing a recipe on the test set, not against a recipe that is simply worse: comparisons of fine-tuned models with API systems
should report how sensitive they are to the training recipe.

\paragraph{Errors are partly complementary.} Decision models and compact classifiers trained on all data share only part of their errors
(Table~\ref{tab:overlap}). This suggests, without testing it, that a combination---for example a compact classifier that defers to a decision
model when it is unsure---might outperform both; the teacher--student and cascade designs discussed in the literature are natural next steps.
